\chapter*{Preface}
\addcontentsline{toc}{chapter}{Preface}
\label{ch:preface}

\section*{About This Manual}

This manual documents SIRCADA, the SIRC Analogue Display Adapter: the picture processing unit (PPU) that turns tile,
tilemap and palette data in its own video memory into a 256 by 224 pixel picture. SIRCADA is a component in its own
right, in a package of its own. It is the companion to the \archmanual{}, which defines the SIRC-1 CPU, to the
\hwmanual{}, which defines the CPU package and bus, and to the \asmmanual{}, which defines the language used to program
them.

SIRCADA is a peripheral on a 16-bit bus, and any 16-bit processor that can supply a chip select, a bus strobe and a
read/write line can drive it. It was designed to work with the SIRC-1, so the signal names, the interrupt arrangements
and the programming examples in this manual are those of the SIRC-1.

It is written for:

\begin{itemize}
	\item Programmers writing graphics code for SIRCADA
	\item Authors of the tools that prepare tiles, tilemaps and palettes
	\item Engineers implementing the device, in hardware or in software
\end{itemize}

\section*{Assumed Knowledge}

This manual assumes familiarity with tile-based two-dimensional graphics hardware: tilemaps, palettes, scanline
rendering and the blanking intervals of a television signal. Readers who need the interrupt mechanism that delivers the
vertical sync signal to the CPU should read the exception chapters of the \archmanual{}.

\section*{Revision Status}

This is a working draft of the \videomanual{}, not for distribution. The device it describes is not finished.
Chapter~\ref{ch:status} lists, feature by feature, what has been built and what has not, and the questions still open;
it is a working record, and is not part of the specification of the device. Anything not yet decided elsewhere in the
manual is marked ``not yet specified''; the manual does not invent values for it.
